\section{Instruments}\label{sec:arch-instruments}

% ─────────────────────────────────────────────────────────────────────────────
\subsection{Products, cash flows and indices}\label{sec:arch-products}
\code{Product} and the instruments; \code{flows/} coupons and legs; generators.

\pending

Equity options derive from \code{Option}, which holds the option type, the strike (a float64 variable), the
exercise type and the underlying. \code{VanillaOption} is European or American; \code{DigitalOption} is European, with a
\code{payoff\_type} (\code{CashOrNothing} or \code{AssetOrNothing}, enum \code{PayoffType}) and a \code{cash\_amount}.
\code{BlackScholesPricer} prices both, the digitals through \code{digital\_calc}; the factory maps a digital to it.
\code{VanillaMCPricer} prices both too; for the digital payoff \code{digital\_smoothing} sets the sigmoid width (default
1\% of the strike, 0 for the hard indicator). \nb{digital\_options.ipynb}

The barrier option class, \code{BarrierOption} in \code{instruments/barrier.py}, adds the level \code{barrier}, the
type \code{barrier\_type} (an enum \code{BarrierType} with the members \code{DownIn}, \code{DownOut}, \code{UpIn} and
\code{UpOut}, and the flags \code{is\_down} and \code{is\_in}) and the dates \code{monitoring\_dates}, where \code{None}
means continuous monitoring (property \code{continuous}). \code{BlackScholesPricer} prices it through \code{barrier\_calc}
(\code{pricers/black.py}); \code{VanillaMCPricer} through \code{\_bridge\_survival} (continuous) or
\code{\_discrete\_survival} (monitoring dates merged into the grid).

With the autocall at an unreachable level, an \code{AutocallableOption} is a strip of digitals (put barrier 0) or a
put knocked in at expiry (coupons 0).
\nb{digital\_strip.ipynb}, \nb{down\_and\_in\_put.ipynb}
